\par\SourceLeaf{152}
It now remains for me to consider a little more closely several points, some of which have already been touched on above.

We have seen above how P. and R. held entirely opposite views on a basic principle (pp. [133] and [137]. When the latter speaks on p. [138, top] of a rise and \textit{fall} of the \textit{fund}, I suspect that he really meant only the \textit{income} from the fund. For it seems to me that, insofar as the institution is a benefaction, it is the strict duty of its management to ensure that the principal from which the benefaction flows is preserved intact. Even the losses on loans of capital that cannot from time to time be avoided despite every precaution make it impossible to ensure this unless, besides preservation, some gradual increase is also made the aim, always preferring to do a little too much rather than too little. In this way the \textit{whole body} of later recipients will not be shortchanged relative to the whole body of earlier recipients, but will rather tend to be better off. This is an advantage that the present recipients can all the more readily wish them to enjoy, since they themselves enjoy the fruits of similar restraint by \textit{their} predecessors. Fairness also requires an effort, through appropriate measures, to offset as far as possible the inevitable fluctuations in the number of recipients, which are proportionately very large in a small society. By contrast, P.'s demand that no future \textit{individual} recipient should ever receive less than any earlier recipient has received seems to me to have no basis in right for a society that, as \textsc{Sch.} very rightly observed, will always retain something of the character of a joint stock company. In any case, and this is the main point, such a demand cannot be met, if it is to count as unconditional, without the most obvious unfairness toward the present recipients. It is evident: the greater the certainty required that this case should never occur, the less may be given to the present recipients. The most extreme cases for the possible number of recipients would have to be considered, from which P.'s assumptions, as shown above, were far removed. This becomes still more striking in the estimates below on p. [150], based on an increased number of participants. Even in 1794 this eventuality could have been included among the future cases that were not merely possible but even probable. For at that time the total number of university professors was 45\footnote{One of them, K., died before the deliberations had ended, on 1794 August 21. Now (in the autumn of 1845) there are 57. But the number of married members of the widows' fund has increased in a much greater proportion, from 26 in 1794 to 42 in 1845.}, and fifty years earlier only 25. It is impossible to determine in advance what limits the continually increasing needs of the times will reach.

The second point concerns the interpretation of the progression rule in \textsc{Sch.}'s sense. When \textsc{Sch.} says [p. 138, line 19] that P. had determined the amount of the increase, and likewise of the reduction, by an unalterable rule, this necessarily presupposes the first interpretation on p. [129], but \textsc{Sch.} nevertheless concedes too much. An unalterable rule for reduction was contained in it only insofar as it specified \textit{when} a reduction must occur (namely, on that interpretation, immediately after the number of 15 pensions was exceeded), but not yet the \textit{size} of the reduction. This incompleteness could, it seems to me, most suitably have been supplied, given the circumstances at that time, by a provision worded as follows:
\begin{quote}
Whenever the capital assets have increased by 5000 Rthl., an increase in pensions takes place, amounting to 10 Rthl. for each individual pension if and for as long as no more than 15 pen%
\SourceLeaf{153}sions exist, but otherwise to a total sum of 150 Rthl., to be divided equally among the individual pensions. The same applies to each subsequent increase.
\end{quote}
\noindent In this way it would not even have been necessary to make the \textit{beginning} of an increase depend on the condition that the number of pensions should not exceed 15. In this respect, this form of provision would even have proved more advantageous to the widows. If a small excess persisted, the loss arising from the obligation to share the previous total increase among more recipients would soon be compensated, or more than compensated, by a new increase, whereas in such a case a new increase is not admissible at all under P.'s rule on the second interpretation.

Moreover, I am inclined to believe that the amended wording above, had anyone thought of proposing it at the time, would also have had to satisfy P. For his assumption that 15 was the highest number of pensions that could occur was either correct or incorrect. In the first case the amendment had no effect at all and was therefore immaterial; in the second it was necessary.

The consequence of such a rule would have been that people would have become accustomed to regarding the pension as composed of two parts: an unchangeable part and an allowance that could from time to time grow with the capital assets, but might also fall back somewhat in the process. The latter, however, would then follow a truly unalterable, simple rule by which each widow could easily check the matter herself. I have already drawn attention above on p. [131] to the distinction between such a reduction governed by a rule and one occurring under the receipt clause owing to the fund's inability to pay.

